% Supplement to A.4.42.  Keep the marking correction before the assertion
% that descent and ordinary forgetting agree on the ordinary normalizer.
\paragraph{The exact marking of ordinary forgetting.}
Write $j_0^{\mathrm{up}}$ for ordinary forgetting with the primitive
upstairs handle basis $(A,B)=(b^3,a^3)$ identified, in that order, with
the generators $(a,b)$ of the target $\Sigma$.  Put
\[
 J=\begin{pmatrix}0&1\\1&0\end{pmatrix},\qquad
 s_J:\Sigma\longrightarrow\Sigma,\quad a\longmapsto b,\ b\longmapsto a.
                                                               \tag{MFP.1}
\]
This is an actual automorphism of the compact target
$\mathcal P_c$.  Indeed $N_8$ is characteristic by (QS.1),
$L_8=\gamma_3(N_8)N_8^3$ and
$T=N_8\cap\ker(\Sigma\to C_3^2)$ are characteristic, and $s_J$
carries the normal closure of $[a,b]^2$ onto itself: it interchanges
the two commutator entries, hence sends their commutator to its inverse.
It also preserves $\Sigma'\Sigma^{12}$.  It therefore preserves the
literal kernel $R_c=(L_8T'D_\partial)\cap\Sigma'\Sigma^{12}$
of (QF.1).  The induced automorphism of $\mathcal P_c$ normalizes
the ordinary image $J_c$ in $\operatorname{Out}(\mathcal P_c)$:
the free-group automorphism $s_J$ is represented by the orientation
reversing interchange of the two handle curves, and conjugation by it
preserves the oriented once-punctured-torus mapping class group.
Consequently there is the specified involution
\[
 \sigma_J:J_c\longrightarrow J_c,\quad
 [\alpha]\longmapsto[s_J\alpha s_J^{-1}],\qquad
 \varpi\bigl(\sigma_J(j)\bigr)=J\varpi(j)J,
                                                               \tag{MFP.2}
\]
where $\varpi:J_c\twoheadrightarrow\operatorname{SL}_2(\mathbb Z/12)$
is the retained signed-period map.  Inner choices of based lifts have
no effect on this formula in the outer image.

In (FCD.1)--(FCD.4) take the aligned ordinary forgetting map to be
\[
                  j_0=\sigma_J\circ j_0^{\mathrm{up}}.
                                                               \tag{MFP.3}
\]
This specifies its target marking without altering any kernel or
identity-coordinate condition used in (FCD.5)--(CDF.9).  The nine induced
representations before and after this choice are related by applying
$\sigma_J$ to each of their nine weights, with the same permutation
coordinate.  Thus their actual joint source images are isomorphic by
this displayed coordinate map.

The homology map of the three-division covering, in the named bases,
is $3J$.  For a zero-fixing actual ordinary deck normalizer its upstairs
and downstairs signed matrices therefore satisfy
\[
       C_{\mathrm{down}}(3J)=(3J)C_{\mathrm{up}},\qquad
       C_{\mathrm{down}}=JC_{\mathrm{up}}J.
                                                               \tag{MFP.4}
\]
The second equality follows in the integral matrix ring from the first,
before reduction modulo twelve.  The oriented once-punctured-torus
mapping class is determined by its integral handle matrix.  Applying
its actual action on $\mathcal P_c$ proves the precise equality
$d\rho=j_0$ on this ordinary normalizer.  With the unaligned convention
the equality instead reads $d\rho=\sigma_J j_0^{\mathrm{up}}$.

\paragraph{The full joint image of the two finite outputs.}
Retain the aligned convention (MFP.3) and put
\[
 \mathcal N_0^\dagger
   =N_{G^\dagger}(H^\dagger)\cap\operatorname{im}(\mathcal M_0).
\]
The exact image of the specified homomorphism on this group is
\[
 \operatorname{im}\!\left(
   (d\operatorname{pr}_G,j_0^\dagger):
   \mathcal N_0^\dagger\longrightarrow J_c\times J_c\right)
     =J_c\mathbin{\times}_{\operatorname{SL}_2(\mathbb Z/12)}J_c.
                                                               \tag{MFP.5}
\]
In particular its order is
\[
                  1152\cdot3^8=7\,558\,272.
                                                               \tag{MFP.6}
\]
To prove the inclusion in (MFP.5), take an actual source representative
in $\mathcal M_0$ of a finite normalizer element.  Its signed upstairs
period $C$ is retained in $G^\dagger$.  The actual finite reconstruction
in (NC.1)--(NC.5) and (TF.4) gives
$\varpi(d\operatorname{pr}_G)=JCJ$.  Ordinary forgetting with (MFP.3)
has this same period.  This argument uses finite normalizer membership
only for the domain of $d$; it does not assert that the source
representative normalizes the deck group before passage to the
finite quotient.

For the reverse inclusion, every integral matrix in
$\operatorname{SL}_2(\mathbb Z)$ is represented by the corresponding
linear map of the torus.  Its lift through multiplication by three
fixes zero and permutes the nine division points; it is an actual
ordinary deck normalizer.  Formula (MFP.4) and the ordinary descent
comparison consequently put the full diagonal $\{(j,j):j\in J_c\}$
in the image of (MFP.5).  The proved equality (CDF.1)--(CDF.9) puts
$\ker\varpi\times\{1\}$ in the image, with source elements even
centralizing $H^\dagger$.  When $\varpi(a)=\varpi(b)$ the literal
factorization
\[
                   (a,b)=(ab^{-1},1)(b,b)
                                                               \tag{MFP.7}
\]
therefore proves surjectivity onto the stated fiber product.
Its order is $|J_c|\,|\ker\varpi|$; A.4.22 and (CDF.9) give (MFP.6).
The same equality holds after adjoining the additional finite source
coordinates with solvable image specified in A.4.42: the kernel
construction retains value one in those coordinates, and actual
ordinary normalizers still supply the full diagonal with their actual
extra values.

Under the unaligned convention the exact image is instead the
explicitly isomorphic twisted fiber product
$\{(a,b):\varpi(a)=J\varpi(b)J\}$, generated by
$\ker\varpi\times\{1\}$ and
$\{(\sigma_J(b),b):b\in J_c\}$.  The isomorphism to (MFP.5) is
$(a,b)\mapsto(a,\sigma_J(b))$.
